% 08. kapitulua - Ariketak: kaixo, Scrumban agentikoa!
% Helburua: zeregin bat, agente bat, pull request bat. Doako mailarekin edo rol-jokoan (ikasleak agente gisa) egin daitekeen lehen ariketa.
% Aurreko liburutik berrerabil daitekeena:
%   - aurreko_liburua_2024/kapituluak/5-prestaketa_lanak.tex (ariketen formatua: "N. ariketa - ...", "Ariketa ebatziak", "Proposatutako ariketa gehigarriak")
% Irudiak: Pictures/08/ (goiburu hutsa header_08.pdf.xoj)
\todo[inline]{Laburpena (behin-behinekoa): Kapitulu honetan lehen aldiz jarriko dugu agente bat lanean gure fluxuan, ahalik eta modurik txikienean: zeregin bat, agente bat, \textit{pull request} bat. \ref{cha:taula}.~kapituluko taula eta \ref{cha:agenteak}.~kapituluko agentea hartuko ditugu, eta pausoz pauso egingo dugu bidea: zeregina idatzi onarpen-irizpideekin, agenteari esleitu, agenteak adarra sortu eta kodea idatzi, \textit{pull request}-a ireki, integrazio jarraituak testak exekutatu, gizakiak berrikusi eta bateratu, eta taulako txartela Eginda zutabera iritsi. Ariketa bakoitzean zerbait apurtuko dugu nahita ---zeregin lausoa, test faltsua, baimen gehiegi--- fluxuak nola erantzuten duen ikusteko. Ariketak doako mailan edo rol-jokoan egin daitezke, ikasleek agentearen papera eginez, tresnarik gabe ere.}

\section{Ariketa ebatziak}

\section{Proposatutako ariketa gehigarriak}
